\clearpage
\section{Receiving, rendering, and becoming willing}
\zlabel{triptych:concise:commentary:start}
\subsection{Does a useful power know whom it serves?}
The troubling words in Isaiah are not only the title ``anointed,'' but the repeated assertion that Cyrus does not know God. Taken alone, the title might suggest a direct identification with Christ. Taken alone, the ignorance might suggest that nothing good could come through the king. Jerome allows neither conclusion. He reads the proper name as Cyrus, identifies the royal office, and lets the negative clause test the proposed identification. The Son who reconciles the world cannot be the unknowing recipient described here. The actual king remains within the literal account.\footnote{Jerome, \work{Commentarii in Isaiam}, XII, on 45:1--7, DCO \work{Pars 3}, pp.126--129, especially 127--128.}

This is also what gives Jerome's typology its precision. A type anticipates a fuller work through a real, limited action. Cyrus's liberation can image Christ's liberation without transferring Cyrus's ignorance to Christ. The fulfilment is greater than the anticipation. Literal distinction and allegorical connection therefore support one another: erase the king, and the concrete act of liberation disappears; equate him with Christ, and the oracle's denial of knowledge becomes unintelligible.

Aquinas examines the benefit from another direction. In Isaiah 45 he distinguishes liberation for the people, recognition of God, and the knowledge extending to the nations. He says Cyrus remained in idolatry. The effects intended through the ruler and the change needed within him cannot be treated as a single accomplishment. The people can receive a genuine benefit even while its instrument fails to acknowledge the giver.\footnote{Aquinas, \work{Expositio super Isaiam ad litteram}, 45, Corpus Thomisticum, paragraph 85005.}

These distinctions yield two different judgments. Ordered allegiance asks how Israel can accept liberation without making its agent a saviour. Willing service asks how the recipient of strength can become a grateful knower of its source. The first guards the use of a good; the second tests its user's heart. Neither permits the conclusion that useful action certifies everything about the actor. Neither requires ingratitude for a benefit simply because the human agent is imperfect.

Jerome's final application makes the inward question unavoidable. He warns a teacher endowed with speech and wisdom to give the glory to God. Otherwise the rebuke to Cyrus can become the teacher's own. Religious subject matter does not secure a religious intention. A person may speak effectively of God while taking the gift as proof of personal independence. That is why the Collect's concern for sincere service reaches farther than competence. It asks for the will that can use a gift truthfully, not merely for another gift to use.\footnote{Jerome, \work{In Isaiam}, XII, conclusion of the exposition of 45:1--7, pp.128--129.}

\subsection{What does the coin establish, and what does it leave open?}
Chrysostom begins with the structure of the trap. An answer against tribute could expose Jesus to Roman power; an answer in its favour could discredit him before the people. The questioners' praise of his truthfulness is therefore an instrument of manipulation. Jesus asks them for the coin and makes them identify its image and inscription. Their answer supplies evidence drawn from the civic exchange in which they already participate. The dilemma no longer belongs entirely to those who devised it.\footnote{Chrysostom, \work{Homilies on Matthew}, LXX §§1--2, NPNF I.10 (1888).}

For Chrysostom, ``render'' means that something is due. It does not describe an arbitrary gift from a person with no obligation. Yet the same exposition refuses an unlimited conclusion: compliance must not injure godliness. Office identifies the person making a demand; office alone does not make every demand righteous. The literal account thus preserves an earthly duty while denying that earthly authority establishes its own final moral measure.

Aquinas follows the literal exchange with an expressly mystical explanation. Caesar's image identifies the coin; God's image identifies the soul that must be returned to him. The two claims are not equal claims on two equal objects. The person who handles the manufactured coin receives being from God. Giving a religious portion of life cannot discharge the whole of that relation, as though all remaining decisions were beyond it. Aquinas also retains the need to seek peace within shared civic life. Belonging to God can generate responsibilities among neighbours.\footnote{Aquinas, \work{Super Matthaeum}, 22, lectio 2, paragraph 87389, concluding literal and mystical explanations.}

Here the two interpretations need one another. An account of ordered duties can remain merely external: the correct amount is paid, while resentment, vanity or dishonesty govern the giver. The interpretation of willing service asks about the person making the return. Conversely, an account of inward sincerity can evade the coin altogether. Aquinas's spiritual extension does not convert an unpaid obligation into a merely interior symbol. The will must become truthful in the practical action it performs.

The Collect and the Prayer over the Offerings bring this distinction into worship. The former concerns readiness and sincere service; the latter joins the servant's offering to purification. A gift can be objectively good while its giver still needs cleansing. That need is not a reason to stop every imperfect good act. It is a reason to receive help within the act, instead of using the act to prove that help is unnecessary. The Gospel's flattery supplies the contrary image: true words become untrue service when they are used to control the one they praise.

\subsection{Why thank God for work that people actually do?}
Paul's three nouns resist passivity. Faith has work, love has labour, and hope has endurance. Chrysostom explains them against the pressure borne by the community: belief remains steadfast, charity spends itself, and hope lasts through affliction. The fuller chapter describes reception amid suffering; that context explains the virtues, although the appointed reading stops at 1:5b. Thanksgiving neither denies the cost nor assigns the whole achievement to the people. Their conduct gives Paul reason to thank God, and thanking God teaches them humility.\footnote{Chrysostom, \work{Homilies on 1 Thessalonians}, I, on 1:1--7, NPNF I.13 (1889). Verse 6 and the final conduct clause of verse 5 are context beyond the Sunday appointment.}

Aquinas asks what the effects reveal. Election has not been earned by prior merit; it is recognized through what occurs in hearer and preacher. His treatment of power, Spirit and fullness allows more than one distribution: signs and spiritual efficacy, or the preacher's virtuous life, the Spirit and complete instruction. The divine efficacy of the Gospel cannot consequently be reduced to a demand for one spectacular sign. Costly fidelity and truthful instruction have their place in the account.\footnote{Aquinas, \work{Super I Thessalonicenses}, 1, lectio 1, paragraph 87862.}

For ordered allegiance, this community supplies a measure of fruitfulness different from dominion. Paul does not thank God because the Thessalonians command the city. Their identity is relational: they are a church in God the Father and the Lord Jesus Christ. The Psalm's call to all peoples can take local shape in a vulnerable community whose hope endures. Visible power and divine election are not interchangeable predicates.

For willing service, the same community answers the apparent competition between gift and effort. A received power can become a genuinely human act. Aquinas develops the point in the context surrounding the acclamation, Philippians 2:12--13: God's inward action reaches willing and doing. He excludes both the claim that free choice saves without help and the claim that providence eliminates free choice. God is not a second human worker who must push the person aside in order to act.\footnote{Aquinas, \work{Super Philippenses}, 2, lectio 3, paragraph 87838. Philippians 2:12--13 is explanatory context, not the acclamation's appointed words.}

Nor can willing be assigned entirely to an independent self, with God responsible only for completing what that self first decided. The passage names willing as well as doing. Prayer for a ready will is therefore meaningful: it asks for help where choosing itself occurs. The moral result is strenuous response without self-origination. Grace is not a reason to postpone charity until it feels effortless; charity's labour is not a price that purchases grace. Paul's thanksgiving holds those truths together in the memory of an actual people.

\subsection{Whose glory becomes visible in the world?}
The responsorial Psalm and acclamation make the inward question public. Augustine understands the Psalm's new song as new love, building a house throughout the world. He then tests the preacher through the object of proclamation: God's glory must be announced, not the speaker's. A successful religious performance can fail at precisely that point. Jerome's gifted teacher and Augustine's preacher face the same temptation from different passages: the received ability becomes an occasion to claim the giver's credit.\footnote{Augustine, \work{Expositions on the Psalms}, XCV §§1--6, NPNF I.8, English heading ``Psalm XCVI.''}

Aquinas's explanation of the acclamation gives visibility a source. God is light by essence; the saints shine by participation in the Word's light. Holding Christ's word and shining are therefore connected acts, rather than competing paths of speech and conduct. The community's visible life need not become invisible in order to confess dependence. Its actual brightness is received brightness.\footnote{Aquinas, \work{Super Philippenses}, 2, lectio 4, paragraph 87839.}

Chrysostom's explanation accents a life-giving principle or seed in the saints. The NPNF editor prefers the Gospel-word sense of ``word of life'' and says so in a note. Aquinas explicitly speaks of Christ's word. These explanations can support a relation between received life and visible conduct without becoming identical accounts of the phrase. Chrysostom is not correcting himself through his editor's note. His surrounding exhortation to serve without murmuring also directs attention to the quality of willing, not simply to public effectiveness.\footnote{Chrysostom, \work{Homilies on Philippians}, VIII, on 2:12--16, including the NPNF editor's note at ``word of life.''}

The Psalm's Christological reception requires another distinction. Augustine's text at XCV §11 includes ``from the wood,'' which he reads of the cross. The appointed Psalm does not contain that phrase, and its strophe omits verse 10b concerning the world's stability. John Paul II expressly presents the cross interpretation as Christian rereading and relates divine kingship to the Son's service in Mark. The allegorical sense can therefore interpret lordship through self-gift while the literal sung text continues to proclaim the Creator's kingship and equitable government.\footnote{Augustine, \work{Expositions}, XCV §11; John Paul II, General Audience, 18 September 2002, §§4--5.}

Ordered allegiance finds here a kingship beyond the claims of any empire. Willing service finds the humility that changes the use of power within the Church. These are complementary emphases, not a quarrel between authors. Their shared moral test is concrete: does a received capacity direct others toward God, or make the servant the final object of admiration?

\subsection{Can dependence give courage rather than diminish it?}
The Entrance's sheltering image gives these public demands a vulnerable speaker. Psalm 16 in Vulgate numbering, 17 in Hebrew numbering, asks for a just hearing, firm steps and protection; its wider conclusion hopes to behold God's glory. The antiphon selects the appeal and shelter, not the whole case for vindication. Fidelity therefore begins neither with possession of public power nor with an assurance that every enemy has disappeared. It begins with a cry addressed to God.

Augustine's explanation of verse 6 notices the order of that cry: the speaker calls because God has heard. Earlier mercy has strengthened prayer that was once weaker. The person does not first produce an independent confidence and then use it to obtain divine attention. Confidence itself already answers help. Augustine places this voice in Christ and his body, allowing the Church to pray from received mercy rather than from a claim to self-sufficient righteousness.\footnote{Augustine, \work{Expositions}, XVI §§1,6,8,14, NPNF I.8, English heading ``Psalm XVII.''}

Aquinas begins with the careful protection suggested by the pupil of the eye. His spiritual applications of the wings include charity and the cross. These are several applications of an image, not one literal description of divine anatomy. Protection under the cross cannot be a promise that Christ's members will escape all suffering. It names a refuge within which suffering can be borne.\footnote{Aquinas, \work{Super Psalmos}, 16 §§3--4, paragraph 86901 and the following division.}

This matters differently to the two interpretations. Ordered allegiance calls for the courage to fulfil a genuine duty and to refuse a demand that violates godliness. That courage cannot rest on the fantasy of invulnerability. Willing service calls for the humility to ask for help even in asking: the will's renewal includes receiving a stronger capacity to pray. Dependence does not abolish either kind of action. It removes the requirement that the actor already possess sufficient strength before responding.

The moral and final horizons remain distinct. Courage concerns a task that must be performed amid present weakness. Hope looks beyond the present to the vision and nourishment developed in the wider Psalm contexts. The final prayer's eternal gifts do not promise that every earthly dispute will end favourably. They make it possible to serve faithfully without requiring earthly success to establish the ultimate worth of the service.

\subsection{What is received before service is offered?}
The Psalm Communion option recalls a Lord whose gaze is merciful. In its full biblical context, armies, kings and horses cannot provide final deliverance. Augustine explains famine as the deprivation of the Word and everlasting truth, then describes patient expectation of incorruptible food. Bodily provisions matter, but they cannot satisfy every hunger or abolish mortality. His account gives ordered allegiance a limit upon the promises of worldly strength and gives willing service a reason to persevere while need remains.\footnote{Augustine, \work{Expositions}, XXXII §§16--22, NPNF I.8, English heading ``Psalm XXXIII.''}

Aquinas differentiates bodily rescue, rescue from sin and final rescue; nourishment likewise receives bodily, spiritual and sacramental explanation. His sacramental reference makes the Communion connection specific. He also joins fear to hope: fear alone tends toward despair, while hope separated from fear becomes presumption. Mercy does not require pretending that sin is harmless or that the needy are already full. It makes truthful reception possible. Neither author promises that faithful people never suffer bodily deprivation.\footnote{Aquinas, \work{Super Psalmos}, 32 §§15--17, especially §16, paragraph 87103.}

The Mark alternative names the prior gift more sharply. Its appointed excerpt is the ransom clause; the preceding service clause and the apostles' dispute belong to the full Gospel context. Bede first describes the Son of Man taking the servant's form and shedding his blood for the world. He then explains ``for many'' as those willing to believe. The gift's worldwide horizon and the believer's response stand together. His gloss does not make faith a payment that precedes redemption, nor does it say that everyone already receives its fruit.\footnote{Bede, \work{In Marci evangelium expositio}, III, Mark 10:42--45, PL 92 (1862), cols.236A--B, especially 236B.}

Bede answers a question about reception; his short gloss does not explain the origin of a will to believe. Aquinas's account of divine action within willing addresses that further question. Chrysostom, commenting on Matthew 20:28 rather than directly on Mark, develops the consequence of the Lord's descent: death for enemies defeats the disciples' contest for status. The disciple receives liberation and then imitates the liberator's humility. He does not become a second redeemer whose service supplies a ransom equal to Christ's.\footnote{Chrysostom, \work{Homilies on Matthew}, LXV §§4--5; Aquinas, \work{Super Philippenses}, 2.3.}

The two Communion options consequently give different shape to the same dependence. The Psalm accents mercy that sustains the needy through fear and hope. Mark accents the saving act received in faith before it becomes a pattern of action. Neither option is an obligatory sequel to the other. With either, the Prayer after Communion directs present benefit toward eternal gifts. A real reception begins growth; it does not announce the end of the journey.

\subsection{How does ransom free rather than merely inspire?}
Ransom makes the Mark option more than an exhortation to helpfulness. Gregory Nazianzen asks to whom the blood could be paid. He rejects a devil entitled to receive God and rejects an account in which the Father demands or delights in blood. He speaks positively of the Father's acceptance, the Incarnation's sanctifying work and liberation through the Son. The saving act cannot be explained by imagining a divine appetite for violence.\footnote{Gregory Nazianzen, \work{Oration} XLV §§21--24, especially §22, Browne--Swallow translation, NPNF II.7 (1894).}

Aquinas uses a different positive account. Christ's voluntary loving gift supplies superabundant satisfaction; the price is paid to God rather than the devil. Humanity never ceased to be subject to God's power, but lost union in charity. Redemption therefore concerns a broken relation, not God's recovery of a world that had escaped his sovereignty. Aquinas's account explains why the gift can restore that union.\footnote{Aquinas, \work{Summa theologiae} III, q.48, a.4 and opening of a.5, English Dominican translation.}

Their shared denial of the devil's entitlement does not erase the difference between Gregory's question about acceptance and Aquinas's language of satisfaction paid to God. Nor does that difference make either account a reduction of redemption to moral example. Both insist on a deed the disciple receives. For ordered allegiance, Christ's liberation prevents any earthly power from claiming the place of final saviour. For willing service, the same priority prevents gratitude from becoming a project of self-redemption.

The four senses meet without collapsing here. Literally, the Son's ransom belongs to his journey toward the Passion, while the coin remains a coin and Cyrus a king. Allegorically, limited liberation and the returning divine image find their fulfilment in Christ and his body. Morally, the worshipper renders true duties and serves through a will receiving help. Anagogically, present nourishment sustains hope for a good not yet fully possessed. The first interpretation frees hope from the final claims of earthly power; the second gives unfinished growth a final direction. The closing prayer joins them: a life truly active can remain a life wholly dependent on gift.
